% siminos/gudorf/ppt/ppt.tex        pdflatex ppt; biber ppt
% $Author: predrag $ $Date: 2019-03-27 16:44:17 -0400 (Wed, 27 Mar 2019) $

%% started with siminos/presentations/MATT18/MATT18.tex
%% started with siminos/gudorf/thesis/proposal/tex/ppt.tex
%% logical setup, no need to edit %%%%%%%%%%
\newif\ifpaper \newif\ifPDF               %%
\newif\ifOUP \newif\ifboyscout            %%
\newif\ifdasbuch \newif\ifarticle         %%
\newif\ifsolutions

\newif\ifbackground % whether use background image
\backgroundfalse
\newif\iflogo % whether use logo
\logotrue
\newif\ifnotes % whether show notes.
\notesfalse
\newif\ifhandout % handout mode
\handoutfalse
%%%%%%%%%%%%%%%%%%%%%%%%%%%%%%%%%%%%%%%%%%%%%%%%%%%%%%%%%%%%%%%%%%%%%%

\ifhandout
\documentclass[mathserif, handout]{beamer}
\else
\documentclass[mathserif]{beamer}
\fi

\usepackage{graphicx}
\usepackage{algpseudocode}
\usepackage{algorithm}
\usepackage{tikz}
\usetikzlibrary{calc}
\usetikzlibrary{decorations.pathreplacing}
\usepackage{tipa}
\usepackage[font=small,labelfont=bf]{caption}
\usepackage{subcaption}
\usepackage{amsthm}

\usepackage[
    backend=biber,  %bibtex,
    sorting=nyt,
    %refsection=chapter,
    %citereset=chapter,
    style=numeric, %alphabetic, % %style=authoryear,
    natbib=true,
    style=phys, % aps
    biblabel= brackets, % superscript, %
    articletitle=false, % true,  % false, % aps
    %chaptertitle=true,  % aip;  % false, % aps
    pageranges = true , % aip: the full range
             % = false, % aps: only the first page being printed
    sortlocale=en_US,
    firstinits=true,
    url=false, %true,  %
    doi=false, %true,
    eprint=false
]{biblatex}
\addbibresource{../../bibtex/siminos.bib}

%%%%%%%%%%%%%%%%%%%%%%%%%%%%%%%%%%%%%%%%%%%%%%%%%%%%%%%%%%%%%%%%%%%%%%
% deal with the conflicts between def.tex and beamer
\newcounter{chapter}
\let\block\undefined
\let\Loop\undefined
\let\definition\undefined
\let\theorem\undefined
\let\lemma\undefined

\input{../../inputs/def}            %% do not edit; update from
\input{../../inputs/editsDasbuch}   %% editing comments, DasBuch style
%\input ../../inputs/defsSpatiotemp
% siminos/gudorf/thesis/defsThesis.tex
% $Author: predrag $ $Date: 2018-11-02 16:08:17 -0400 (Fri, 02 Nov 2018) $

\ifboyscout %%%%%%%% DISPLAY COMMENTS IN THE TEXT %%%%%%%%%%%%%%%%%%%%
            %%%%%%%% turn on labeling of equations on margins %%%%%%%%
    % also search the text for lines starting with %%  to
    % locate various internal comments, recent edits etc.
    \typeout{============ COMMENTED =====}
 \renewcommand{\PC}[2]{$\footnotemark\footnotetext{Predrag #1: #2}$} %date, comment
 \newcommand{\AC}[2]{$\footnotemark\footnotetext{Andy #1: #2}$} %date, comment
 \newcommand{\BG}[2]{$\footnotemark\footnotetext{Boris #1: #2}$} %date, comment
 \renewcommand{\AKS}[2]{$\footnotemark\footnotetext{Adrien #1: #2}$} %date, comment
 \renewcommand{\RJ}[2]{$\footnotemark\footnotetext{Rana #1: #2}$} %date, comment
 \renewcommand{\HL}[2]{$\footnotemark\footnotetext{Han #1: #2}$} %date, comment
 \newcommand{\MNG}[2]{$\footnotemark\footnotetext{Matt #1: #2}$} %date, comment
\else % drop comments
      % do not turn on labeling of equations on margins
  \typeout{============ UNCOMMENTED =====}
 \renewcommand{\PC}[2]{}
 \newcommand{\AC}[2]{}
 \newcommand{\BG}[2]{}
 \renewcommand{\AKS}[2]{}
 \renewcommand{\RJ}[2]{}
 \renewcommand{\HL}[2]{}
 \newcommand{\MNG}[2]{}
\fi  %%%%%%%%%%%% END OF ON/OFF COMMENTS SWITCH %%%%%%%%%%%%%%%%%%%%

%\newcommand{\YTlink}[1]{\marginpar{   %% links to YouTube videos
%    \HREF{http://#1}{\includegraphics[width=16pt]{start-32}}}
%                       }
\newcommand{\YTlink}[1]{    %% in text, links to YouTube videos
    \raisebox{-0.4ex}[0pt][0pt]{\!\!\!
    \HREF{http://#1}{\includegraphics[height=1em]{start-32}}
    \!\!\!            }  }
% in-text pointer to a video   2014-01-18
\newcommand{\videoLink}[1]{    %% in text, links to YouTube videos
    \raisebox{-0.4ex}[0pt][0pt]{$\!\!\!$
    \HREF{http://#1}{\includegraphics[height=1em]{YouTubeIcon}}
                    % or maybe streamThisRed}}
                    % do not have compact eps for {start-32}
    $\!$             }}
% in-text pointer to a video   2014-01-18
\newcommand{\toVideo}[1]
{\marginpar[\videoLink{#1}]{\videoLink{#1} \hfill}}

% \renewcommand{\toChap}[1]{}
\renewcommand{\cl}[1]{{\ensuremath{|#1|}}}  % the length of a periodic orbit, Ronnie
%\newcommand{\XD}{Xiong Ding}

\ifpaper % prepare for B&W paper printing:
       \renewcommand{\arXiv}[1]{ {arXiv:#1}}
\else % prepare hyperlinked pdf
       \renewcommand{\arXiv}[1]{
              {\href{http://arXiv.org/abs/#1}{arXiv:#1}}}
\fi

    \ifboyscout
\newcommand{\toCB}{\marginpar{\footnotesize 2CB}}  % to compare with ChaosBook
\newcommand{\inCB}{\marginpar{\footnotesize now in CB}} % entered in ChaosBook
\newcommand{\CBlibrary}[1]
{\href{http://ChaosBook.org/library/#1.pdf} { (click here)}}
    \else
\newcommand{\toCB}{}
\newcommand{\inCB}{}
\newcommand{\CBlibrary}[1]{}
    \fi %end of internal draft switch

\newcommand{\conf}{\ensuremath{x}}      % Configuration space coordinate
\newcommand{\Fu}{\tilde{u}}
\newcommand{\abs}[1]{\lvert#1\rvert}    % Absolute value
\newcommand{\twot}{invariant 2-torus} % Predrag, to distinguish from a PO
\newcommand{\twots}{invariant 2-tori}
\newcommand{\twoT}{Invariant 2-torus}
\newcommand{\twoTs}{Invariant 2-tori}
\newcommand{\brick}{block}
\newcommand{\Brick}{Block}
\renewcommand{\action}{W}                 % PC 07jan2018
% \newcommand{\action}{S}               % Boris
%\newcommand{\genF}{F}                  % Li and Tomsovic generating function
\newcommand{\genF}{L}                   % McKay generating function

%%%%%% Boris definitions
\newcommand{\Aa}{\ensuremath{\bar{\A}}}
\newcommand{\A}{\ensuremath{\mathcal{A}}}       % alphabet
\newcommand{\Ai}{\ensuremath{\mathcal{A}_0}}    % alphabet inner
\newcommand{\Ae}{\ensuremath{\mathcal{A}_1}}    % alphabet outer
\newcommand{\R}{\ensuremath{\mathcal{R}}}
%\newcommand{\m}{\ensuremath{\mathsf{m}}}     % Boris
\newcommand{\m}{\ensuremath{m}}     % Predrag experimental 2016-11-08
\newcommand{\Mm}{\ensuremath{\mathsf{M}}}
\newcommand{\Xx}{\ensuremath{\mathsf{X}}}
\newcommand{\Zz}{\ensuremath{\mathbb{Z}^2}}
\newcommand{\ZLT}{\mathbb{Z}^2_{\scriptscriptstyle\mathrm{LT}}}
%\newcommand{\q}{\ensuremath{\mathsf{m}}}     % Boris
\newcommand{\q}{\ensuremath{q}}     % Predrag experimental 2016-11-08
\newcommand{\D}{\mathcal{D}}
\newcommand{\gd}{\mathsf{g}}
\newcommand{\gp}{\mathsf{g}^0}

%%%%%%%%%%%% MACROS, Xiong Ding specific %%%%%%%%%%
\DeclareMathOperator{\im}{Im}
\DeclareMathOperator{\re}{Re}
\DeclareMathOperator{\diag}{Diag}
\DeclareMathOperator{\dist}{dist}
\DeclareMathOperator{\const}{const}
% ==============================
\newcommand{\spt}{spatiotemporal}
\newcommand{\Spt}{Spatiotemporal}
\newcommand{\xDFt}{Discrete Fourier transform}
\newcommand{\xdFt}{discrete Fourier transform}

\newcommand{\cGL}{complex Ginzburg-Landau}
\newcommand{\cGLe}{complex Ginzburg-Landau equation}


%%%%%%%%%%%% MACROS, Siminos specific %%%%%%%%%%
\newcommand{\edit}[1]{{\color{blue} #1}} % for referees
%\newcommand{\edit}[1]{#1}               % for the journal
\newcommand{\vf}{v}	%%% keep notation for vector field flexible.
                    % For the time being follow Das Buch.
\newcommand{\conj}[1]{\ensuremath{\bar{#1}}}
\newcommand{\Manif}{\ensuremath{\mathcal{M}}}
\newcommand{\Order}[1]{\mathrm{O}(#1)}
\newcommand{\steady}{\textdollar~}

%%%%%%%% Symmetries, Siminos specific %%%%%%%%%%
%\renewcommand{\shift}{\ensuremath{\ell}}
%\newcommand{\Shift}{\ensuremath{\tau}}
\newcommand{\Idg}{\ensuremath{\mathbf{1}}}
\newcommand{\En}[1]{\ensuremath{\mathrm{E}(#1)}}
\newcommand{\Rotn}[1]{\ensuremath{R_{#1}}}
\newcommand{\Drot}{\ensuremath{\zeta}}
\newcommand{\globstab}[1]{\ensuremath{\Sigma_{#1}}} % Change to be the same as stab. Was \Sigma^\ast_{#1}
\newcommand{\Str}[1]{\ensuremath{\mathcal{S}_{#1}}} % Stratum
\newcommand{\Nlz}[1]{\ensuremath{N(#1)}}
\newcommand{\doubleperiod}[1]{{\ensuremath{\mathcal{T}_{#1}}}}

%%%%%%%%%%%% Theorems, Siminos specific %%%%%%%%%%
%%%%%%%%%%%%%% Solution labels %%%%%%%%%%%%%%%%%%%%


% Redefine using mathrm, it is a label not a math symbol
\renewcommand{\EQV}[1]{\ensuremath{\mathrm{E}_{#1}}}
% E_0: u = 0 - trivial equilibrium
% E_1,E_2,E_3, for 1,2,3-wave equilibria
\renewcommand{\REQV}[2]{\ensuremath{\mathrm{TW}_{#1#2}}} % #1 is + or -
% TW_1^{+,-} for 1-wave traveling waves (positive and negative velocity).
\renewcommand{\PO}[1]{\ensuremath{\mathrm{PO}_{#1}}}
% PO_{period to 2-4 significant digits} - periodic orbits
\renewcommand{\RPO}[1]{\ensuremath{\mathrm{RPO}_{#1}}}
% RPO_{period to 2-4 significant digits} - relative PO.  We use ^{+,-}
% to distinguish between members of a reflection-symmetric pair.
% Gibson likes:
\renewcommand{\tEQ}{\ensuremath{\mathrm{EQ}}}

\newcommand{\tLB}{\ensuremath{{\text{EQ}_1}}}
\newcommand{\tUB}{\ensuremath{{\text{EQ}_2}}}

%%%%%%%%%%%%%% Operators %%%%%%%%%%%%%%%%%%%%%%%
\newcommand{\PperpOp}{\mathbf{P}^{\perp}}
%\newcommand{\Pperp}{P^{\perp}}
%

%%%%%%%%%%% Equations, Siminos specific %%%%%%%%%%
\newcommand{\cont}{\,, \\ }

%%%%%%%%%%%%%%%%%%%%%%%%%%%%%%%%%%%%%%%%%%%%%%%%%%%%%%%%%%%
% JFG favorite macros
\newcommand{\bu}{\ensuremath{{\bf u}}}
\newcommand{\bx}{{\bf x}}
\newcommand{\trHalf}[1]{\tau_{#1}}    % 1/2 cell translation
\newcommand{\GPCF}{\ensuremath{\Gamma}} % Hoyle notation, equivariant symmetry group

%%%%%%%%%%%%%%%%%%%%%%%%%%%%%%%%%%%%%%%%%%%%%%%%%%%%%%%%%%%
% GUDORF specific macros.
\newcommand{\F}{\ensuremath{\mathcal{F}}}
\newcommand{\Fi}{\ensuremath{\mathcal{F}^{-1}}}
\newcommand{\Dxx}{\ensuremath{\mathcal{D}_{\conf \conf}}}
\newcommand{\Dfourx}{\ensuremath{\mathcal{D}_{\conf \conf \conf \conf}}}
\newcommand{\Sop}{\ensuremath{\mathcal{S}}}
\newcommand{\Dt}{\ensuremath{\mathcal{D}_{\zeit}}}
\newcommand{\Dx}{\ensuremath{\mathcal{D}_{\conf}}}
\newcommand{\ufield}{\ensuremath{u(\conf,\zeit)}}
\newcommand{\dufield}{\ensuremath{u(\conf_{m},\zeit_{n})}}
\newcommand{\sFb}{spatiotemporal Fourier basis}
\newcommand{\rv}{real valued}
\newcommand{\Fuvec}{\ensuremath{\mathbf{\tilde{u}}}}
\newcommand{\umn}{\ensuremath{\tilde{u}_{\sss{m,n}}}}
\newcommand{\nm}[1]{\tilde{{#1}}_{\sss{n,m}}}
\newcommand{\qm}{\ensuremath{q_{\sss m}}}
\newcommand{\wavevec}{\ensuremath{\frac{2\pi m}{L}}}
\newcommand{\wn}{\ensuremath{\omega_{\sss n}}}
\newcommand{\freqvec}{\ensuremath{\frac{2\pi n}{T}}}

%%%%%%%%%%%%%%%%%%%%%%%%%%%%%%%%%%%%%%%%%%%%%%%%%%%%%%%%%%%
% Lopez favorite macros, REMOVE MOST EVENTUALLY
\newcommand{\sss}{\scriptscriptstyle}	
\newcommand{\e}{\mathrm{e}}
\newcommand{\mmbf}[1]{\mbox{\boldmath ${#1}$}}
\newcommand{\Arot}{\theta}
\newcommand{\spacetrans}{\sigma}
\newcommand{\timetrans}{\tau}
\newcommand{\lsym}{l}
\newcommand{\ce}{c_{\sss{1}}}
\newcommand{\co}{c_{\sss{2}}}
\newcommand{\Lx}{L_{\sss x}}
\newcommand{\omegan}{\omega_{\sss n}}
\newcommand{\km}{k_{\sss m}}
\newcommand{\am}[1]{a_{\sss{{#1}}}}
\newcommand{\amn}[2]{\hat{a}_{\sss{{#1},{#2}}}}
\newcommand{\bmn}[2]{\hat{b}_{\sss{{#1},{#2}}}}
\newcommand{\torus}{{\mathbb T}}
\newcommand{\dd}{\mathrm{d}}
\newcommand{\cglegroup}{G}
\newcommand{\cgleisotropy}[1]{\cglegroup_{\!{\sss{{#1}}}}}
\newcommand{\baseSpace}{\mathcal{B}}
\newcommand{\moduliSpace}{\pS} %\EuFrak{M}}
% \newcommand{\moduliSpaceC}{\mathcal{M}} % conflicts with \pS
\newcommand{\moduliSpaceI}{\mathcal{I}}
\newcommand{\solSpace}{\mathcal{S}}
\newcommand{\fiberM}[2]{\mathcal{#1}_{\sss{#2}}}


%%%%%%%%%%%% REMOVE THIS EVENTUALLY %%%%%%%%%%%
%


\graphicspath{{../../figs/ksstFigs/}{../../figs/}{../../Fig/}}

% PC temporary fix, will not be needed in later MixTeX version         2018-11-02
% https://tex.stackexchange.com/questions/426088/texlive-pretest-2018-beamer-and-subfig-collide
\makeatletter
\let\@@magyar@captionfix\relax
\makeatother

% ------                       setup of beamer               ----------------
\usetheme{Madrid}

\ifnotes
\setbeameroption{show notes}
\else
\setbeameroption{hide notes}
\fi

\setbeamertemplate{bibliography item}[text]
% add background image
\ifbackground
\usebackgroundtemplate
{
  \includegraphics[width=\paperwidth,height=\paperheight]{./figs/background.jpg}%
}
\fi

% -------------------      title page       --------------------------
\title[Spatiotemporal K-S equation ] % (optional, only for long titles)
{}
\subtitle{Spatiotemporal Tiling of the \KSe\, }
\logo{
  \iflogo
  \includegraphics[width=2cm,height=1cm,keepaspectratio]{gtlogo}
  \includegraphics[width=1cm,height=0.8cm,keepaspectratio]{cnslogo}
  \fi
}
\author[M. Gudorf] % (optional, for multiple authors)
{Matthew Gudorf %\inst{1}
}
\institute[GAtech] % (optional)
{
  % \inst{1}%
  Center for Nonlinear Science \\
  School of Physics \\
  Georgia Institute of Technology
}

\date[January 25th, 2019] % (optional)
{ Adviser: Predrag Cvitanovi\'c \\ \vspace{1em}
} %\\ Topics in Geometric and Nonlinear Dynamical Systems}

\subject{Physics}

% ----------------    setup of outline    -------------------------
\AtBeginSection[]
{
  \begin{frame}
    \frametitle{Table of Contents}
    \tableofcontents[currentsection]
  \end{frame}
}

% ---------------   main content of slides  -----------------------
\begin{document}
\frame{\titlepage} % add title page.

% ----------------  section: definition of problem --------------


%\begin{frame}
%  \frametitle{Spatiotemporal Chaos}
%\begin{itemize}
%\item What do I mean by spatiotemporal chaos?
%Fully developed chaos; takes up full domain.
%Large number of length scales.
%Deterministic, (usually)hyperbolic,ergodic systems
%\item Why study spatiotemporal chaos?
%Many macroscopic models of natural processes exhibit spatiotemporal chaos.
%Infinite limit contains infinite amount of scales

%\item What will I focus on?
%Many different aspects; pattern formation, dynamical averages, analytic
%solutions for some systems, intermittency, statistical calculations,
%weak chaos near nonequilibrium pattern formation is very common
%Time invariant solutions, spatiotemporal recurrences.
%\end{itemize}
%\end{frame}

\begin{frame}
\frametitle{Motivation}
\begin{itemize}
\item High dimensional fluid flows are hard. Most results
in ``minimal flow units''\footfullcite{JM91}
\item Conventional computational tools are providing diminishing returns
\item Spatiotemporal methods are scalable\footfullcite{WGBGQ13}
\end{itemize}

\end{frame}



\begin{frame}
  \frametitle{ Doubly Periodic \KSe\ }
\begin{figure}[h]
\includegraphics[width=.8\textwidth]{MNG_annotatedks}
\end{figure}
Boundary conditions $u(x,t)=u(x,t+T_p)=u(x+L,t)=u(x+L,t+T_p)$
%Nice problem to work with: two dimensional-> easy to visualize, easier than Navier-Stokes,
%models wide variety of phenomena
%Large space or time makes things hard
%Recurrences are spatiotemporal
%Relatively small number of patterns appear frequently
\end{frame}


\begin{frame}
    \frametitle{Large spatiotemporal simulation}
\begin{figure}[h]
\includegraphics[width=.6\textwidth]{MNG_uu500b500}
\end{figure}
\end{frame}

\begin{frame}
    \frametitle{Shadowing}
\begin{figure}[h]
\includegraphics[width=.48\textwidth]{MNG_uu500b500co}
\includegraphics[width=.48\textwidth]{cutouts}
\end{figure}
\end{frame}
%Talk about the motivation to formulate spatiotemporal symbolic dynamics
%here.

%Main purpose of this work is to develop new techniques through a spatiotemporal formulation
%to find solutions on larger and larger domains. Then use symbolic dynamics to


\begin{frame}
\frametitle{Fourier-Fourier Basis}
%Motivated by smooth solutions, diagonal basis for translations.
    Substituting
    \beq
    u(x,t) = \sum_{m,n} \umn e^{\ii \wn t + \ii \qm x}
    \,,
    \eeq
    into the \KSe\ gives,

    \beq \label{eq:Fks}
    \mathbf{F}(\mathbf{x})= (\ii \freqvec + (\wavevec)^2 - (\wavevec)^4 ) \umn + \frac{1}{2} \wavevec \mathcal{F}((\mathcal{F}^{-1}(\umn))^2)
    \eeq
    Solve nonlinear algebraic equations for $\mathbf{x} = (\Fuvec,T,L)$.
    Similar to formulation of Lopez\footfullcite{Lopez2015}.
    We use real-valued Fourier basis in computations.

\end{frame}

\begin{frame}
  \frametitle{Disadvantages of spatiotemporal methods}
  \begin{enumerate}
  \item High dimensionality
  \item Numerically challenging
  \item Possible to find isolated solutions that are unimportant
  \item Find different members of same continuous families (redundancy)
  \end{enumerate}
\end{frame}

\begin{frame}
  \frametitle{Advantages of spatiotemporal methods}
  \begin{enumerate}
  \item Can use variational methods
  \item Fourier-Fourier basis
  \item Topological constraint (boundary conditions) provides
  access to potent methods
  \item Stability does not play a role in finding solutions.
  \item Can exploit spatiotemporal symmetries
  \end{enumerate}
\end{frame}


\begin{frame}
  \frametitle{Spatiotemporal Symmetries}
  Reflection invariance
  \beq \nonumber
  \ufield = \sigma_{\conf}  \ufield
  \eeq
  \begin{figure}[h]
  \includegraphics[width=.9\textwidth]{REFL}
  \end{figure}
\end{frame}

\begin{frame}
  \frametitle{Spatiotemporal Symmetries}
  Shift-Reflection Invariance
  \beq \nonumber
  \ufield = \sigma_{\conf} \tau_{1/2} \ufield
  \eeq
  \begin{figure}[h]
  \includegraphics[width=.9\textwidth]{SR}
  \end{figure}
\end{frame}

\begin{frame}
  \frametitle{Spatiotemporal Symmetries}
  Spatial Translation Symmetry
  \[
  \ufield = g(\theta(t)) \circ u(x,t+T_p)
  \]
  \begin{figure}[h]
  \includegraphics[width=.6\textwidth]{RPO}
  \end{figure}
\end{frame}

\begin{frame}%[allowframebreaks]
  \frametitle{Initial conditions}
  \begin{figure}
  \includegraphics[width=.32\textwidth]{RPOinit}
  \includegraphics[width=.32\textwidth]{SRinit}
  \includegraphics[width=.32\textwidth]{ANTIinit}
  \end{figure}
\end{frame}

\begin{frame}
  \frametitle{Adjoint Descent Method}
    \beq
    \mathbf{F}(\mathbf{x})= (\ii \freqvec + (\wavevec)^2 - (\wavevec)^4 ) \umn + \frac{1}{2} \wavevec \mathcal{F}((\mathcal{F}^{-1}(\umn))^2)
    \eeq
  \beq
  I = \frac{1}{2} \mathbf{F}^{\top}\mathbf{F}
  \,.
  \eeq
\end{frame}

\begin{frame}
  \frametitle{Adjoint Descent Method}
  introduce fictitious time ($\tau$) flow by differentiation of cost function,
  \beq
  \partial_{\tau}I = (J^{\top}\mathbf{F})^{\top}(\partial_{\tau}\mathbf{x})
  \eeq
  ``adjoint descent'' method is defined by choice\footfullcite{Faraz15}
  \beq
  \partial_{\tau}\mathbf{x} = -(J^{\top}\mathbf{F})
  \eeq
\end{frame}

%This can be expensive to evaluate so in reality the code uses a version
%that calculates this quantity without ever forming the matrix J.

\begin{frame}
  \frametitle{Damped Gauss-Newton Method}
  Solve $\mathbf{F}(\mathbf{x}^*)=0$ via linearized system
  \beq
  \mathbf{F}(\mathbf{x}^*) = \mathbf{F}(\mathbf{x}) + J\delta \mathbf{x}
  \eeq
  \beq
  J = \begin{bmatrix}
  \frac{\partial \mathbf{F}}{\partial \mathbf{x}} &
  \frac{\partial \mathbf{F}}{\partial T} &
  \frac{\partial \mathbf{F}}{\partial L}
  \end{bmatrix}
  \eeq
  \beq
  J \delta \mathbf{x} = -\mathbf{F}(\mathbf{x})
  \eeq
\end{frame}


\begin{frame}%[allowframebreaks]
  \frametitle{Shift-reflect symmetry}
  A couple thousand new \twot\ solutions, examples:
  \begin{figure}
  \includegraphics[width=.9\textwidth]{PPO_IF}
  \end{figure}
\end{frame}

\begin{frame}%[allowframebreaks]
  \frametitle{Spatial translation symmetry}
  \begin{figure}
  \includegraphics[width=.9\textwidth]{RPO_IF}
  \end{figure}
\end{frame}

\begin{frame}%[allowframebreaks]
  \frametitle{Reflection symmetry}
  \begin{figure}
  \includegraphics[width=.9\textwidth]{ANTI_IF}
  \end{figure}
\end{frame}

\begin{frame}%[allowframebreaks]
  \frametitle{Qualitative guesses for tiles}
  Sift through collection to find frequently occurring \spt\ patterns
  \begin{figure}
  \includegraphics[width=.2\textwidth]{MNG_defect}
  \includegraphics[width=.2\textwidth]{MNG_hook}
  \includegraphics[width=.2\textwidth]{MNG_half}
  \includegraphics[width=.2\textwidth]{MNG_gap}
  \includegraphics[width=.2\textwidth]{MNG_streak}
  \end{figure}
\end{frame}

\begin{frame}%[allowframebreaks]
  \frametitle{Qualitative tiling of actual solution}
  \begin{figure}
    \includegraphics[width=.4\textwidth,height=.4\textheight]{MNG_ppo_L30_T44}
    \includegraphics[width=.28\textwidth,height=.35\textheight]{MNG_ppo_frankenstein}
  \end{figure}
\end{frame}

\begin{frame}%[allowframebreaks]
  \frametitle{Finding tiles by converging subdomains}
  \begin{figure}
  \includegraphics[width=.2\textwidth]{gap0}
  \includegraphics[width=.2\textwidth]{gap1}
  \includegraphics[width=.2\textwidth]{gap2}
  \includegraphics[width=.2\textwidth]{MNG_gap_final}
  \end{figure}
\end{frame}

\begin{frame}%[allowframebreaks]
  \frametitle{Converged tiles}
  \begin{figure}
  \includegraphics[width=.2\textwidth]{MNG_defect_final}
  \includegraphics[width=.2\textwidth]{MNG_hook_final}
  \includegraphics[width=.2\textwidth]{MNG_gap_final}
  \includegraphics[width=.2\textwidth]{MNG_streak_final}
  \end{figure}
\end{frame}

\begin{frame}%[allowframebreaks]
  \frametitle{Continuous families and numerical continuation}
  \begin{figure}
  \includegraphics[width=.4\textwidth]{MNG_hookdefect_comparison}
  \end{figure}
\end{frame}


\begin{frame}%[allowframebreaks]
  \frametitle{ How do we use tiles?}
  \begin{figure}
  \includegraphics[width=.22\textwidth]{MNG_half_final}
  \includegraphics[width=.1\textwidth,height=.25\textheight]{MNG11symb}
  \includegraphics[width=.22\textwidth]{MNG11init}
  \includegraphics[width=.22\textwidth]{MNG11final}
  \end{figure}
\end{frame}

\begin{frame}%[allowframebreaks]
  \frametitle{Spatiotemporal symbolic dynamics}
\begin{figure}
\includegraphics[width=.25\textwidth,height=.25\textheight]{MNG_HOD_init}
\includegraphics[width=.25\textwidth,height=.25\textheight]{MNG_HOD_final}
\includegraphics[width=.15\textwidth,height=.25\textheight]{MNG_symb11rot}
\end{figure}
\end{frame}

\begin{frame}%[allowframebreaks]
  \frametitle{Spatiotemporal symbolic dynamics}
\begin{figure}
\includegraphics[width=.25\textwidth,height=.25\textheight]{MNG_HOD_STREAK_init}
\includegraphics[width=.25\textwidth,height=.25\textheight]{MNG_HOD_STREAK_final}
\includegraphics[width=.25\textwidth,height=.25\textheight]{MNG_symb1100rot}
\end{figure}
\end{frame}


\begin{frame}%[allowframebreaks]
  \frametitle{Spatiotemporal symbolic dynamics}
\begin{figure}
\includegraphics[width=.25\textwidth,height=.25\textheight]{MNG_HD_init}
\includegraphics[width=.25\textwidth,height=.25\textheight]{MNG_HD_STREAK_init}
\includegraphics[width=.25\textwidth,height=.25\textheight]{MNG_HD_STREAK_final}
\includegraphics[width=.25\textwidth,height=.25\textheight]{MNG_symb1222rot}
\end{figure}
\end{frame}

\begin{frame}%[allowframebreaks]
  \frametitle{Quantitative tiling of actual solution}
  \begin{figure}
  \includegraphics[width=.25\textwidth,height=.3\textheight]{MNG_ppo_frankenstein}
  \includegraphics[width=.25\textwidth,height=.3\textheight]{MNG_approxsymbdyn}
  \includegraphics[width=.25\textwidth,height=.3\textheight]{MNG_approxsymbdyntrinary}
  \includegraphics[width=.3\textwidth,height=.3\textheight]{MNG_ppo_L30_T44}
  \end{figure}
  \footnotetext{Each symbol is technically paired with spatiotemporal parameters $(T,L)$}
\end{frame}

\begin{frame}
\frametitle{New solutions from gluing together old}
\begin{figure}
\includegraphics[width=.3\textwidth,height=.4\textheight]{MNG_ppo_without_HOD}
\includegraphics[width=.3\textwidth,height=.4\textheight]{MNG_ppo_subdomain_HOD}
\end{figure}
\end{frame}

\begin{frame}
\frametitle{New solutions from gluing together old}
\begin{figure}
\includegraphics[width=.3\textwidth,height=.4\textheight]{MNG_ppo_without_HalfD}
\includegraphics[width=.3\textwidth,height=.4\textheight]{MNG_ppo_subdomain_HalfD}
\end{figure}
\end{frame}

\begin{frame}
\frametitle{New solutions from gluing together old}
\begin{figure}
\includegraphics[width=.3\textwidth,height=.2\textheight]{MNG_ppo_subdomain2_1} \\
\includegraphics[width=.3\textwidth,height=.2\textheight]{MNG_ppo_subdomain1_1} \\
\includegraphics[width=.3\textwidth,height=.2\textheight]{MNG_ppo_subdomain0_1}
\end{figure}
\end{frame}

\begin{frame}
\frametitle{New solutions from gluing together old}
\begin{figure}
\includegraphics[width=.3\textwidth,height=.3\textheight]{MNG_ppo_tiling_init_pretruncation_1}
\includegraphics[width=.3\textwidth,height=.3\textheight]{MNG_ppo_tiling_init_1}
\end{figure}
\end{frame}

\begin{frame}
\frametitle{New solutions from gluing together old}
\begin{figure}
\includegraphics[width=.3\textwidth,height=.3\textheight]{MNG_ppo_tiling_final_1}
\includegraphics[width=.3\textwidth,height=.3\textheight]{MNG_ppo_L30_T44}
\end{figure}
\end{frame}

\begin{frame}%[allowframebreaks]
  \frametitle{Gluing (larger) solutions, \ppo\ time example}
  \begin{figure}
  \includegraphics[width=.9\textwidth]{MNG_ppo1ppo2_time}
  \end{figure}
\end{frame}

\begin{frame}%[allowframebreaks]
  \frametitle{Gluing (larger) solutions, \rpo\ time example}
  \begin{figure}
  \includegraphics[width=.9\textwidth]{MNG_rpo1rpo2_time}
  \end{figure}
\end{frame}

\begin{frame}%[allowframebreaks]
  \frametitle{Gluing (larger) solutions, \ppo\ space example}
  \begin{figure}
  \includegraphics[width=.6\textwidth]{MNG_ppo1ppo2_space}
  \end{figure}
\end{frame}

\begin{frame}
\frametitle{Future Work}
\begin{itemize}

\item Investigate Lagrangian formulation and application of Hill's formula%
\footfullcite{IbragiKols04}${}^,$%
\footfullcite{BolTre10}
%To determine which solutions are important
\item Complete spatiotemporal symbolic dynamics
%To be able to enumerate all solutions
\item Continuing improvement of numerical methods
%Of course this process never ends
\item Application of theory to PDEs with more continuous spatial dimensions such as
quasi 2-D Kolmogorov flow%
\footfullcite{STGS18}${}^,$%
\footfullcite{ChaKer12}.

\end{itemize}
\end{frame}

\begin{frame}
\frametitle{Conclusion}
\begin{itemize}
\item Developed spatiotemporal methods for finding \twots\ of the \KSe\
\item Beginning to develop symbolic dynamics for the \spt\ \KSe\
\item Developed new method for finding larger (domain) solutions by combining
smaller (domain) solutions
\end{itemize}
\end{frame}

%\begin{frame}[allowframebreaks]
%  \frametitle{References}
%  \scriptsize{\bibliographystyle{apalike}}
%  \bibliography{../../bibtex/siminos}
%\end{frame}




\end{document}
